% Regression: parallel routes on one turned hull edge are model-proven
% disjoint, so the crossing police must not ask PGF's singular line
% intersection arithmetic to weigh them.
\documentclass{standalone}
\usepackage{tenkz}
\makeatletter
\ExplSyntaxOn
\file_input:n { tenkz-kernel.code.tex }
\tl_new:N \l__tenkz_test_route_name_tl
\cs_new_eq:NN \__tenkz_test_string_count:nnN \__tenkz_string_count:nnN
\cs_new_eq:NN
  \__tenkz_test_route_measure:n
  \__tenkz_kernel_r_route_measure:n
\cs_set_protected:Npn \__tenkz_kernel_r_route_measure:n #1
  {
    \__tenkz_test_route_measure:n {#1}
    \__tenkz_model_get:nnN {#1} {name} \l__tenkz_test_route_name_tl
    \str_if_eq:VnT \l__tenkz_test_route_name_tl {endpoint-route}
      {
        \prop_map_inline:Nn \l__tenkz_kernel_ro_family_prop
          {
            \clist_if_in:nnT {##2} {endpoint-route}
              {
                \errmessage
                  {address-bearing~route~was~classified~as~fully~disjoint}
              }
          }
      }
  }
\cs_set_protected:Npn \__tenkz_string_count:nnN #1#2#3
  {
    \str_if_eq:nnT {#1} {lane-a}
      {
        \str_if_eq:nnT {#2} {lane-b}
          { \errmessage{parallel~lanes~were~sent~to~the~intersection~engine} }
      }
    \str_if_eq:nnT {#1} {lane-b}
      {
        \str_if_eq:nnT {#2} {lane-a}
          { \errmessage{parallel~lanes~were~sent~to~the~intersection~engine} }
      }
    \str_if_eq:nnT {#1} {ring-a}
      {
        \str_if_eq:nnT {#2} {ring-b}
          { \errmessage{closed~rings~were~sent~to~the~intersection~engine} }
      }
    \str_if_eq:nnT {#1} {ring-b}
      {
        \str_if_eq:nnT {#2} {ring-a}
          { \errmessage{closed~rings~were~sent~to~the~intersection~engine} }
      }
    \__tenkz_test_string_count:nnN {#1} {#2} #3
  }
\ExplSyntaxOff
\makeatother
\begin{document}
\tenkzkernel
\begin{tenkz}[rows={ket}, cols=6, frame=circle, physical=up]
  \tn[skin=box, name=A]{A} & \tn[skin=box, name=B]{B}
  & \tn[skin=box, name=C]{C} & \tn[skin=box, name=D]{D}
  & \tn[skin=box, name=E]{E} & \tn[skin=box, name=F]{F}
  \tnwire[name=lane-a, kind=string, species=g, crossing=over,
          route={n of (1,1) .. (1,2)}]{open}{open}
  \tnwire[name=lane-b, kind=string, species=h, crossing=under,
          route={n of (1,1) .. (1,2)}]{open}{open}
\end{tenkz}
\quad
\begin{tenkz}[rows={wire}, cols=6, frame=circle, bonds=none]
  \tn[skin=box, name=G]{G} & \tn[skin=box, name=H]{H}
  & \tn[skin=box]{I} & \tn[skin=box]{J}
  & \tn[skin=box]{K} & \tn[skin=box]{L}
  \tnwire[name=endpoint-route, kind=string,
          route={n of (1,1) .. (1,2)}]{G}{H}
\end{tenkz}
\quad
\begin{tenkz}[rows={wire}, cols=4, bonds=none]
  \tn[skin=box]{M} & \tn[skin=box]{N}
  & \tn[skin=box]{O} & \tn[skin=box]{Q}
  \tnwire[name=ring-a, kind=string, closed,
          route={all of (1,1) .. (1,4)}]
  \tnwire[name=ring-b, kind=string, closed,
          route={all of (1,1) .. (1,4)}]
\end{tenkz}
\end{document}
